% ============================================================
%  RELATÓRIO — ARQUITETURAS DE ALTO DESEMPENHO
%  Práticas 1 e 2: Cálculo de π por Monte Carlo e Leibniz
%  Normas para Elaboração de Relatórios ver 4 (2026)
%  Universidade Federal de São Carlos (UFSCar)
% ============================================================
\documentclass[12pt, a4paper]{article}

% --- Codificação e idioma ---
\usepackage[utf8]{inputenc}
\usepackage[T1]{fontenc}
\usepackage{babel}
\usepackage{amsmath}
\usepackage{amssymb}
\usepackage{multirow}

% --- Fonte Times New Roman ---
\usepackage{mathptmx}

% --- Geometria da página ---
\usepackage[
top=2.5cm,
bottom=2.5cm,
left=3.0cm,
right=3.0cm
]{geometry}

% --- Espaçamento entre linhas ---
\usepackage{setspace}
\onehalfspacing

% --- Indentação do parágrafo ---
\usepackage{indentfirst}
\setlength{\parindent}{1.25cm}

% --- Pacotes gráficos e de figuras ---
\usepackage{graphicx}
\usepackage{float}

% --- Tabelas ---
\usepackage{booktabs}
\usepackage{array}
\usepackage{longtable}

% --- Listagens de código ---
\usepackage{listings}
\usepackage{xcolor}

\lstdefinestyle{matlab}{
	language=Matlab,
	basicstyle=\ttfamily\small,
	keywordstyle=\color{blue}\bfseries,
	commentstyle=\color{gray}\itshape,
	stringstyle=\color{teal},
	numbers=left,
	numberstyle=\tiny\color{gray},
	stepnumber=1,
	numbersep=8pt,
	frame=single,
	breaklines=true,           
	breakatwhitespace=true,   
	columns=flexible,          
	keepspaces=true,          
	captionpos=b,
	morekeywords={parfor,tic,toc,format,rand,sum,pi,length,sqrt}
}

\lstdefinestyle{python}{
	language=Python,
	basicstyle=\ttfamily\small,
	keywordstyle=\color{blue}\bfseries,
	commentstyle=\color{gray}\itshape,
	stringstyle=\color{teal},
	numbers=left,
	numberstyle=\tiny\color{gray},
	stepnumber=1,
	numbersep=8pt,
	frame=single,
	breaklines=true,
	breakatwhitespace=true,   
	columns=flexible,          
	keepspaces=true,          
	captionpos=b,
	morekeywords={np,multiprocessing,Pool,time}
}

% --- Numeração de seções ---
\setcounter{secnumdepth}{3}

% --- Cabeçalho e rodapé ---
\usepackage{fancyhdr}
\setlength{\headheight}{14pt}
\pagestyle{fancy}
\fancyhf{}
\fancyhead[R]{\small Arquiteturas de Alto Desempenho --- Relatório}
\fancyfoot[C]{\thepage}
\renewcommand{\headrulewidth}{0.4pt}

% --- Hiperlinks (DEVE ser o último pacote) ---
\usepackage[hidelinks]{hyperref}

% ============================================================
%  INÍCIO DO DOCUMENTO
% ============================================================
\begin{document}
	
	% --- Sem numeração de página na capa ---
	\thispagestyle{empty}
	
	% ============================================================
	%  CAPA
	% ============================================================
	\begin{titlepage}
		\centering
		
		\noindent\rule{\textwidth}{0.8pt}\\[0.2cm]
		{\large \textbf{UNIVERSIDADE FEDERAL DE SÃO CARLOS}}\\[0.15cm]
		{\normalsize Centro de Ciências Exatas e Tecnologia}\\[0.15cm]
		{\normalsize Departamento de Computação}\\[0.15cm]
		
		\vspace{0pt plus 0.35fill}
		
		{\large \textbf{Laboratório de Arquiteturas de Alto Desempenho}}\\[0.4cm]
		{\large Práticas 1 e 2 — Cálculo de $\pi$ por Monte Carlo e pela Série de Leibniz}\\[1.0cm]
		{\normalsize Professor: \textit{Dr. Emerson Carlos Pedrino}}\\
		
		\vspace{0pt plus 0.65fill}
		
		\begin{center}
			{\normalsize \textit{Pedro Augusto Faria} \ (RA: \textit{821124})}\\[0.15cm]

			{\normalsize Curso: \textit{Engenharia da Computação}}\\[0.15cm]
		\end{center}
		
	\end{titlepage}
	
	\setcounter{page}{1}
	
	% ============================================================
	%  SEÇÃO 1 — INTRODUÇÃO                                [1,0 pt]
	% ============================================================
	\section{Introdução}
	
	O cálculo do número $\pi$ constitui um problema clássico que, ao longo da história
	da computação, tem servido como \textit{benchmark} natural para avaliar tanto a
	precisão de métodos numéricos quanto o desempenho de arquiteturas de processamento.
	Este relatório documenta a realização das Práticas~1 e~2 da disciplina
	\textbf{Arquiteturas de Alto Desempenho} (1001537 -- UFSCar), cujo objetivo central
	é comparar duas abordagens algorítmicas para a estimação de $\pi$ -- o Método de
	Monte Carlo e a Série de Leibniz -- sob três estratégias de implementação distintas:
	vetorizada, serial e paralela.
	
	A relevância do estudo reside menos na estimativa de $\pi$ em si e mais na
	exploração dos fundamentos práticos do processamento de alto desempenho: de que
	modo a estrutura de um algoritmo determina seu potencial de paralelização? Quando
	a vetorização é vantajosa e quando se torna um gargalo? Como o \textit{overhead}
	de paralelismo condiciona o ganho real de desempenho? Essas questões permeiam toda
	a análise apresentada nas seções seguintes.
	
	\subsection{Fundamentação Teórica}
	
	\subsubsection{Aritmética de Ponto Flutuante e Representação de $\pi$}
	
	O valor exato de $\pi$ é irracional e transcendente, portanto, sua representação
	computacional é sempre uma aproximação de precisão finita. Em MATLAB/Octave, a
	constante \texttt{pi} é armazenada como um \texttt{double} de 64~bits (padrão
	IEEE~754), com cerca de 15--16 dígitos decimais significativos. O erro de
	aproximação de qualquer método numérico é medido em relação a esse valor de
	referência:
	
	\begin{equation}
		\varepsilon = \left| \pi_{\text{aprox}} - \pi_{\text{ref}} \right|
		\label{eq:erro}
	\end{equation}
	
	\subsubsection{Speedup e Lei de Amdahl}
	
	O \textit{speedup} quantifica o ganho de desempenho de uma implementação em relação
	a uma linha de base serial:
	
	\begin{equation}
		S = \frac{T_{\text{serial}}}{T_{\text{paralelo}}}
		\label{eq:speedup}
	\end{equation}
	
	A Lei de Amdahl estabelece o limite teórico desse ganho em função da fração
	paralelizável $f$ do programa e do número de processadores $p$:
	
	\begin{equation}
		S_{\max} = \frac{1}{(1-f) + \dfrac{f}{p}}
		\label{eq:amdahl}
	\end{equation}
	
	\noindent Quando o \textit{overhead} de paralelização (distribuição de tarefas,
	sincronização e comunicação de memória) é considerado, o speedup observado tende
	a ser inferior ao previsto pela Equação~\eqref{eq:amdahl}, podendo inclusive ser
	menor que~1 para cargas de trabalho muito granulares.
	
	\subsection{Objetivos}
	
	As práticas têm os seguintes objetivos:
	
	\begin{itemize}
		\item Implementar o cálculo de $\pi$ pelo Método de Monte Carlo nas versões
		vetorizada, serial e paralela em MATLAB;
		\item Implementar o cálculo de $\pi$ pela Série de Leibniz nas mesmas três
		versões;
		\item Medir e comparar tempos de execução, erros de aproximação e
		\textit{speedups} entre as implementações;
		\item Analisar o comportamento de desempenho em função do número de amostras
		$n$, variando de valores pequenos a $n = 5\times10^7$;
		\item Identificar a arquitetura paralela mais adequada a cada método, com base
		nos resultados obtidos e nos conceitos teóricos da disciplina;
		\item Comparar os resultados do MATLAB com implementações equivalentes em Python
		(biblioteca \texttt{multiprocessing} na Prática~1 e \texttt{numpy} na
		Prática~2).
	\end{itemize}
	
	\subsection{Metodologia}
	
	Cada método de estimação de $\pi$ foi implementado em três versões independentes
	no MATLAB: \textbf{(i)} vetorizada, explorando operações sobre arrays sem laços
	explícitos; \textbf{(ii)} serial, com laço \texttt{for} convencional como linha
	de base; e \textbf{(iii)} paralela, utilizando \texttt{parfor} para distribuição
	automática das iterações entre núcleos disponíveis. Os tempos de execução foram
	medidos com \texttt{tic}/\texttt{toc} e o erro calculado conforme a
	Equação~\eqref{eq:erro}. O valor de referência \texttt{pi} foi acessado com
	\texttt{format long} ativado para máxima precisão. Os experimentos foram
	conduzidos variando $n$ e os resultados compilados em tabelas e gráficos de
	\textit{speedup} e erro.
	
	% ============================================================
	%  SEÇÃO 2 — DESENVOLVIMENTO TEÓRICO                   [2,0 pt]
	% ============================================================
	\section{Desenvolvimento Teórico}
	
	\subsection{Método de Monte Carlo para Estimação de $\pi$}
	
	\subsubsection{Dedução Geométrica}
	
	Considera-se um círculo de raio $r = 1$ inscrito num quadrado de lado $2r = 2$,
	ambos centrados na origem do plano cartesiano. A área do círculo e a área do
	quadrado são, respectivamente:
	
	\begin{equation}
		A_{\text{círculo}} = \pi r^2 = \pi
		\qquad \text{e} \qquad
		A_{\text{quadrado}} = (2r)^2 = 4
		\label{eq:areas}
	\end{equation}
	
	A razão entre as duas áreas fornece:
	
	\begin{equation}
		\frac{A_{\text{círculo}}}{A_{\text{quadrado}}} = \frac{\pi}{4}
		\label{eq:razao}
	\end{equation}
	
	O Método de Monte Carlo explora essa razão geometricamente: geram-se $n$ pontos
	$(x_i, y_i)$ com coordenadas sorteadas de forma uniformemente aleatória no
	intervalo $[-1,\,1]$. Um ponto pertence ao interior do círculo se e somente se:
	
	\begin{equation}
		x_i^2 + y_i^2 \leq 1
		\label{eq:circulo}
	\end{equation}
	
	Denominando $n_c$ o número de pontos que satisfazem a condição~\eqref{eq:circulo},
	a fração $n_c/n$ converge para a razão de áreas~\eqref{eq:razao} pela Lei dos
	Grandes Números, resultando na estimativa:
	
	\begin{equation}
		\boxed{\pi \approx 4 \cdot \frac{n_c}{n}}
		\label{eq:montecarlo}
	\end{equation}
	
	\subsubsection{Convergência e Erro}
	
	Por ser um método estocástico, o erro de Monte Carlo decresce segundo o desvio
	padrão amostral, que escala como:
	
	\begin{equation}
		\varepsilon_{\text{MC}} \sim \frac{1}{\sqrt{n}}
		\label{eq:erro_mc}
	\end{equation}
	
	\noindent Isso implica que, para reduzir o erro por um fator de 10, é necessário
	aumentar $n$ por um fator de 100 — convergência lenta, mas o método é
	naturalmente paralelizável pois cada ponto é amostrado de forma independente.
	
	\subsection{Série de Leibniz para $\pi$}
	
	\subsubsection{Dedução via Série de Taylor}
	
	Parte-se da função $\arctan(x)$, cuja expansão em série de Taylor em torno de
	$x = 0$ é:
	
	\begin{equation}
		\arctan(x) = \sum_{k=0}^{\infty} \frac{(-1)^k \, x^{2k+1}}{2k+1}
		= x - \frac{x^3}{3} + \frac{x^5}{5} - \frac{x^7}{7} + \cdots
		\label{eq:taylor_arctan}
	\end{equation}
	
	\noindent convergente para $|x| \leq 1$. Avaliando em $x = 1$ e utilizando
	$\arctan(1) = \pi/4$:
	
	\begin{equation}
		\frac{\pi}{4} = \sum_{k=0}^{\infty} \frac{(-1)^k}{2k+1}
		= 1 - \frac{1}{3} + \frac{1}{5} - \frac{1}{7} + \cdots
		\label{eq:leibniz_quarto}
	\end{equation}
	
	Multiplicando ambos os lados por 4, obtém-se a \textbf{Fórmula de Leibniz}:
	
	\begin{equation}
		\boxed{\pi = 4 \sum_{k=0}^{n} \frac{(-1)^k}{2k+1}}
		\label{eq:leibniz}
	\end{equation}
	
	\noindent onde a igualdade exata é atingida apenas no limite $n \to \infty$.
	
	\subsubsection{Convergência e Erro}
	
	A série de Leibniz é uma série alternante de termos com módulo decrescente a zero.
	Pelo Critério de Leibniz (Teorema das Séries Alternantes), o erro de truncamento
	após $n+1$ termos é limitado pelo módulo do primeiro termo desprezado:
	
	\begin{equation}
		\varepsilon_{\text{Leibniz}} \leq \frac{4}{2(n+1)+1} = \frac{4}{2n+3}
		\approx \frac{2}{n}
		\label{eq:erro_leibniz}
	\end{equation}
	
	\noindent O erro decresce, portanto, como $\mathcal{O}(1/n)$, o que é
	\textbf{mais lento} do que a intuição inicial sugere: embora a taxa formal seja
	$1/n$ para ambos os métodos (compare com~\eqref{eq:erro_mc}, que é $1/\sqrt{n}$),
	na prática a constante multiplicativa do erro de Leibniz é muito menor, resultando
	em erros absolutos muito inferiores para o mesmo $n$ — conforme confirmado pelos
	resultados experimentais ($\varepsilon_\text{Leibniz} \approx 10^{-8}$ contra
	$\varepsilon_\text{MC} \approx 10^{-4}$ para $n = 5 \times 10^7$).
	
	\subsection{Arquitetura Paralela Adequada}
	
	Com base na análise teórica das características de cada método, é possível
	antecipar o tipo de arquitetura paralela mais adequada a cada um.
	
	O Método de Monte Carlo é um exemplo arquetípico de problema \textbf{embaraçosamente
		paralelo} (\textit{embarrassingly parallel}): as amostras são independentes entre
	si, não há comunicação entre tarefas e o custo computacional por iteração é
	suficientemente alto (geração de dois números aleatórios e uma operação de
	comparação) para amortizar o \textit{overhead} de paralelização. Arquiteturas do
	tipo \textbf{SIMD} (Single Instruction, Multiple Data) ou clusters de memória
	distribuída (\textbf{MIMD}) são adequadas, com excelente escalabilidade.
	
	A Série de Leibniz, por sua vez, possui iterações extremamente leves (uma divisão
	e uma adição), tornando o custo de \textit{overhead} de paralelização dominante
	sobre o custo de computação. Para esse tipo de problema, a vetorização via operações
	em bloco (\textbf{SIMD}/\textit{loop unrolling}) é preferível ao paralelismo
	baseado em threads, conforme previsto pela Lei de Amdahl~\eqref{eq:amdahl} e
	corroborado pelos resultados práticos.
	
	% ============================================================
	%  SEÇÃO 3 — DESENVOLVIMENTO PRÁTICO                   [3,0 pt]
	% ============================================================
	\section{Desenvolvimento Prático}
	\label{sec:pratico}
	
	\subsection{Ambiente de Execução}
	
	Foi utilizado um sistema computacional Linux, com especificações a seguir. Além disso, o daemon CPU-X foi usado para obter informações detalhadas das especificações de hardware.
	
	Os experimentos foram realizados integralmente em MATLAB. O ambiente de hardware
	e software utilizado é descrito na Tabela~\ref{tab:ambiente}. As medições de tempo
	empregaram exclusivamente as funções nativas \texttt{tic}/\texttt{toc}, e o valor
	de referência de $\pi$ foi obtido com \texttt{format long} ativado. Além disso, foi utilizado ambiente de desenvolvimento Python, através da IDE VSCode.
	
	\begin{table}[H]
		\centering
		\caption{Especificações do ambiente de execução.}
		\label{tab:ambiente}
		\begin{tabular}{ll}
			\toprule
			\textbf{Parâmetro}          & \textbf{Valor}                          \\
			\midrule
			Software principal          & MATLAB (2019)               \\
			Linguagem comparativa       & Python 3.11 (\texttt{multiprocessing}, \texttt{numpy}) \\
			Processador                 & Intel(R) Core(TM) i5-7400 CPU @ 3.00GHz         \\
			Número de núcleos/threads   & 4/4                            \\
			Memória RAM                 & 16GB                               \\
			Sistema Operacional         & Ubuntu 24.04.4 LTS                               \\
			\bottomrule
		\end{tabular}
	\end{table}
	
	\begin{figure}[H]
		\centering
		\includegraphics[width=1.10\textwidth]{matlab-environment.png}
		\caption{Ambiente MATLAB utilizado durante os experimentos.}
		\label{fig:matlab_env}
	\end{figure}
	
	\begin{figure}[H]
		\centering
		\includegraphics[width=0.80\textwidth]{cpux.png}
		\caption{Interface do CPU-X exibindo as especificações de hardware do sistema.}
		\label{fig:cpux}
	\end{figure}
	
	\subsection{Complete Code Files}
	
	\subsubsection{Implementação de Monte Carlo em MATLAB}
	A implementação em MATLAB para Monte Carlo serial, vetorizado e paralelo.
	
	\begin{lstlisting}[style=matlab, caption={pratica01.m}, label={lst:matlab}]
		% Pedro Augusto Faria - 821124
		% Arquiteturas de Alto Desempenho - Prof. Emerson
		
		% Pratica 01 - Simulacao de Monte Carlo
		
		clear; clc;
		format long;
		
		n = 50000000;
		
		tic;
		
		% 1. Versao Vetorial
		
		x = 2*rand(n,1) - 1;
		y = 2*rand(n,1) - 1;
		
		inside=  (x.^2 + y.^2) <= 1;
		
		pi_vec = 4 * sum(inside) / n;
		
		time_vec = toc;
		
		tic;
		
		% 2. Versao Serial
		
		tic;
		
		count = 0;
		
		for i = 1:n
		x = 2*rand - 1;
		y = 2*rand - 1;
		
		if (x^2 + y^2) <= 1
		count = count + 1;
		end
		end
		
		pi_serial = 4 * count / n;
		
		time_serial = toc;
		
		% 3. Versao Paralela
		
		tic;
		
		count_par = 0;
		
		parfor i = 1:n
		x = 2*rand - 1;
		y = 2*rand - 1;
		
		if (x^2 + y^2) <= 1
		count_par = count_par + 1;
		end
		end
		
		pi_par = 4 * count_par / n;
		
		time_par = toc;
		
		% 4. Calculo de Erro
		
		pi_true = pi;
		
		error_vec = abs(pi_vec - pi_true);
		error_serial = abs(pi_serial - pi_true);
		error_par = abs(pi_par - pi_true);
		
		% 5. Speedup
		
		speedup_vec = time_serial / time_vec;
		speedup_par = time_serial / time_par;
		
		% 6. Display de Resultados
		
		disp('==============================================')
		disp('        MONTE CARLO PI ESTIMATION')
		disp('==============================================')
		
		fprintf('\nNumber of points (n): %d\n', n)
		
		disp('----------------------------------------------')
		disp('Estimated Values of Pi:')
		fprintf('Vectorized  : %.10f\n', pi_vec)
		fprintf('Serial      : %.10f\n', pi_serial)
		fprintf('Parallel    : %.10f\n', pi_par)
		fprintf('MATLAB Pi   : %.10f\n', pi)
		
		disp('----------------------------------------------')
		disp('Absolute Errors:')
		fprintf('Vectorized  : %.10e\n', error_vec)
		fprintf('Serial      : %.10e\n', error_serial)
		fprintf('Parallel    : %.10e\n', error_par)
		
		disp('----------------------------------------------')
		disp('Execution Times (seconds):')
		fprintf('Vectorized  : %.6f s\n', time_vec)
		fprintf('Serial      : %.6f s\n', time_serial)
		fprintf('Parallel    : %.6f s\n', time_par)
		
		disp('----------------------------------------------')
		disp('Speedup (relative to Serial):')
		fprintf('Vectorized  : %.2f x faster\n', speedup_vec)
		fprintf('Parallel    : %.2f x faster\n', speedup_par)
		
		disp('==============================================')
		
		% ==============================================
		% 7. ANALISE DE PERFORMANCE vs n
		% ==============================================
		
		n_values = [1e3, 5e3, 1e4, 5e4, 1e5, 5e5];
		
		time_vec_arr = zeros(size(n_values));
		time_serial_arr = zeros(size(n_values));
		time_par_arr = zeros(size(n_values));
		
		for k = 1:length(n_values)
		
		n = n_values(k);
		
		% -------- Vectorized --------
		tic;
		x = 2*rand(n,1) - 1;
		y = 2*rand(n,1) - 1;
		inside = (x.^2 + y.^2) <= 1;
		pi_vec = 4 * sum(inside) / n;
		time_vec_arr(k) = toc;
		
		% -------- Serial --------
		tic;
		count = 0;
		for i = 1:n
		x = 2*rand - 1;
		y = 2*rand - 1;
		if (x^2 + y^2) <= 1
		count = count + 1;
		end
		end
		time_serial_arr(k) = toc;
		
		% -------- Parallel --------
		tic;
		count_par = 0;
		parfor i = 1:n
		x = 2*rand - 1;
		y = 2*rand - 1;
		if (x^2 + y^2) <= 1
		count_par = count_par + 1;
		end
		end
		time_par_arr(k) = toc;
		end
		
		figure;
		
		plot(n_values, time_vec_arr, '-o', 'LineWidth', 2);
		hold on;
		plot(n_values, time_serial_arr, '-s', 'LineWidth', 2);
		plot(n_values, time_par_arr, '-^', 'LineWidth', 2);
		
		xlabel('Number of points (n)');
		ylabel('Execution Time (seconds)');
		title('Performance Comparison');
		
		legend('Vectorized', 'Serial', 'Parallel');
		grid on;
		
		speedup_vec_arr = time_serial_arr ./ time_vec_arr;
		speedup_par_arr = time_serial_arr ./ time_par_arr;
		
		figure;
		
		plot(n_values, speedup_vec_arr, '-o', 'LineWidth', 2);
		hold on;
		plot(n_values, speedup_par_arr, '-s', 'LineWidth', 2);
		
		xlabel('Number of points (n)');
		ylabel('Speedup');
		title('Speedup vs Problem Size');
		
		legend('Vectorized Speedup', 'Parallel Speedup');
		grid on;
		
		% FIM
	\end{lstlisting}
	
	\subsubsection{Implementação para Leibniz em MATLAB}
	Equivalentemente, a implementação de Leibniz serial, vetorizado e paralelo.
	
	\begin{lstlisting}[style=matlab, caption={pratica02.m}, label={lst:matlab}]
		
	% Pedro Augusto Faria - 821124
	% Arquiteturas de Alto Desempenho - Prof. Emerson
	
	% Pratica 02 - Serie de Leibniz para calculo de Pi
	
	clear; clc;
	format long;
	
	n = 50000000;
	
	% ==============================================
	% 1. VERSAO VETORIZADA
	% ==============================================
	
	tic;
	
	k = (0:n-1)';
	terms = (-1).^k ./ (2*k + 1);
	pi_vec = 4 * sum(terms);
	
	time_vec = toc;
	
	% ==============================================
	% 2. VERSAO SERIAL
	% ==============================================
	
	tic;
	
	sum_serial = 0;
	
	for k = 0:n-1
	sum_serial = sum_serial + ((-1)^k)/(2*k + 1);
	end
	
	pi_serial = 4 * sum_serial;
	
	time_serial = toc;
	
	% ==============================================
	% 3. VERSAO PARALELA
	% ==============================================
	
	
	tic;
	
	sum_par = 0;
	
	parfor k = 0:n-1
	sum_par = sum_par + ((-1)^k)/(2*k + 1);
	end
	
	pi_par = 4 * sum_par;
	
	time_par = toc;
	
	% ==============================================
	% 4. CALCULO DE ERRO
	% ==============================================
	
	pi_true = pi;
	
	error_vec = abs(pi_vec - pi_true);
	error_serial = abs(pi_serial - pi_true);
	error_par = abs(pi_par - pi_true);
	
	% ==============================================
	% 5. SPEEDUP
	% ==============================================
	
	speedup_vec = time_serial / time_vec;
	speedup_par = time_serial / time_par;
	
	% ==============================================
	% 6. DISPLAY DE RESULTADOS
	% ==============================================
	
	disp('==============================================')
	disp('        LEIBNIZ PI ESTIMATION')
	disp('==============================================')
	
	fprintf('\nNumber of terms (n): %d\n', n)
	
	disp('----------------------------------------------')
	disp('Estimated Values of Pi:')
	fprintf('Vectorized  : %.10f\n', pi_vec)
	fprintf('Serial      : %.10f\n', pi_serial)
	fprintf('Parallel    : %.10f\n', pi_par)
	fprintf('MATLAB Pi   : %.10f\n', pi)
	
	disp('----------------------------------------------')
	disp('Absolute Errors:')
	fprintf('Vectorized  : %.10e\n', error_vec)
	fprintf('Serial      : %.10e\n', error_serial)
	fprintf('Parallel    : %.10e\n', error_par)
	
	disp('----------------------------------------------')
	disp('Execution Times (seconds):')
	fprintf('Vectorized  : %.6f s\n', time_vec)
	fprintf('Serial      : %.6f s\n', time_serial)
	fprintf('Parallel    : %.6f s\n', time_par)
	
	disp('----------------------------------------------')
	disp('Speedup (relative to Serial):')
	fprintf('Vectorized  : %.2f x faster\n', speedup_vec)
	fprintf('Parallel    : %.2f x faster\n', speedup_par)
	
	disp('==============================================')
	
	% ==============================================
	% 7. ANALISE DE PERFORMANCE vs n
	% ==============================================
	
	n_values = [1e3, 5e3, 1e4, 5e4, 1e5, 5e5];
	
	time_vec_arr = zeros(size(n_values));
	time_serial_arr = zeros(size(n_values));
	time_par_arr = zeros(size(n_values));
	
	for idx = 1:length(n_values)
	
	n = n_values(idx);
	
	% -------- Vectorized --------
	tic;
	k = (0:n-1)';
	terms = (-1).^k ./ (2*k + 1);
	pi_vec = 4 * sum(terms);
	time_vec_arr(idx) = toc;
	
	% -------- Serial --------
	tic;
	sum_serial = 0;
	for k = 0:n-1
	sum_serial = sum_serial + ((-1)^k)/(2*k + 1);
	end
	time_serial_arr(idx) = toc;
	
	% -------- Parallel --------
	tic;
	sum_par = 0;
	parfor k = 0:n-1
	sum_par = sum_par + ((-1)^k)/(2*k + 1);
	end
	time_par_arr(idx) = toc;
	end
	
	% ==============================================
	% 8. PLOT - TEMPO DE EXECUCAO
	% ==============================================
	
	figure;
	
	plot(n_values, time_vec_arr, '-o', 'LineWidth', 2);
	hold on;
	plot(n_values, time_serial_arr, '-s', 'LineWidth', 2);
	plot(n_values, time_par_arr, '-^', 'LineWidth', 2);
	
	xlabel('Number of terms (n)');
	ylabel('Execution Time (seconds)');
	title('Leibniz Performance Comparison');
	
	legend('Vectorized', 'Serial', 'Parallel');
	grid on;
	
	% ==============================================
	% 9. PLOT - SPEEDUP
	% ==============================================
	
	speedup_vec_arr = time_serial_arr ./ time_vec_arr;
	speedup_par_arr = time_serial_arr ./ time_par_arr;
	
	figure;
	
	plot(n_values, speedup_vec_arr, '-o', 'LineWidth', 2);
	hold on;
	plot(n_values, speedup_par_arr, '-s', 'LineWidth', 2);
	
	xlabel('Number of terms (n)');
	ylabel('Speedup');
	title('Leibniz Speedup vs Problem Size');
	
	legend('Vectorized Speedup', 'Parallel Speedup');
	grid on;
	
	% FIM
		
	\end{lstlisting}
	
	\subsubsection{Implementação Equivalente em Python}
	
	Os dois algoritmos são implementados de forma equivalente em Python, usando as bibliotecas apropriadas.
	
	\begin{lstlisting}[style=python, caption={pratica01e02.py}, label={lst:python_mc}]
		# ============================================================
		# Pedro Augusto Faria - 821124
		# High Performance Architectures
		# Python Benchmark: Monte Carlo + Leibniz
		# ============================================================
		
		import numpy as np
		import time
		import math
		import multiprocessing as mp
		import matplotlib.pyplot as plt
		
		# ============================================================
		# CONFIG
		# ============================================================
		
		N_MONTE = int(5e6)     
		N_LEIBNIZ = int(5e5)
		
		N_VALUES = [int(x) for x in [1e3, 5e3, 1e4, 5e4, 1e5, 5e5]]
		
		# ============================================================
		# MONTE CARLO 
		# ============================================================
		
		def monte_carlo_vectorized(n):
		x = np.random.uniform(-1, 1, n)
		y = np.random.uniform(-1, 1, n)
		inside = (x**2 + y**2) <= 1
		return 4 * np.sum(inside) / n
		
		def monte_carlo_serial(n):
		count = 0
		for _ in range(n):
		x = np.random.uniform(-1, 1)
		y = np.random.uniform(-1, 1)
		if x*x + y*y <= 1:
		count += 1
		return 4 * count / n
		
		def mc_worker(n):
		x = np.random.uniform(-1, 1, n)
		y = np.random.uniform(-1, 1, n)
		return np.sum((x**2 + y**2) <= 1)
		
		def monte_carlo_parallel(n):
		p = mp.cpu_count()
		chunk = n // p
		with mp.Pool(p) as pool:
		results = pool.map(mc_worker, [chunk]*p)
		return 4 * sum(results) / (chunk * p)
		
		# ============================================================
		# LEIBNIZ 
		# ============================================================
		
		def leibniz_vectorized(n):
		k = np.arange(n)
		return 4 * np.sum((-1)**k / (2*k + 1))
		
		def leibniz_serial(n):
		s = 0.0
		for k in range(n):
		s += ((-1)**k) / (2*k + 1)
		return 4 * s
		
		def leibniz_worker(chunk):
		start, end = chunk
		s = 0.0
		for k in range(start, end):
		s += ((-1)**k) / (2*k + 1)
		return s
		
		def leibniz_parallel(n):
		p = mp.cpu_count()
		chunk_size = n // p
		chunks = [(i*chunk_size, (i+1)*chunk_size) for i in range(p)]
		
		with mp.Pool(p) as pool:
		results = pool.map(leibniz_worker, chunks)
		
		return 4 * sum(results)
		
		# ============================================================
		# BENCHMARK FUNCTION
		# ============================================================
		
		def benchmark(method_name, vec_fn, ser_fn, par_fn, n):
		
		print("\n" + "="*50)
		print(f"        {method_name.upper()} PI ESTIMATION")
		print("="*50)
		
		# -------- Vectorized --------
		t0 = time.time()
		pi_vec = vec_fn(n)
		t_vec = time.time() - t0
		
		# -------- Serial --------
		t0 = time.time()
		pi_ser = ser_fn(n)
		t_ser = time.time() - t0
		
		# -------- Parallel --------
		t0 = time.time()
		pi_par = par_fn(n)
		t_par = time.time() - t0
		
		# -------- Errors --------
		pi_true = math.pi
		err_vec = abs(pi_vec - pi_true)
		err_ser = abs(pi_ser - pi_true)
		err_par = abs(pi_par - pi_true)
		
		# -------- Speedup --------
		speed_vec = t_ser / t_vec
		speed_par = t_ser / t_par
		
		# -------- Display --------
		print(f"\nNumber of samples (n): {n}")
		
		print("\nEstimated Pi:")
		print(f"Vectorized  : {pi_vec:.10f}")
		print(f"Serial      : {pi_ser:.10f}")
		print(f"Parallel    : {pi_par:.10f}")
		print(f"Reference   : {pi_true:.10f}")
		
		print("\nAbsolute Errors:")
		print(f"Vectorized  : {err_vec:.2e}")
		print(f"Serial      : {err_ser:.2e}")
		print(f"Parallel    : {err_par:.2e}")
		
		print("\nExecution Time (s):")
		print(f"Vectorized  : {t_vec:.6f}")
		print(f"Serial      : {t_ser:.6f}")
		print(f"Parallel    : {t_par:.6f}")
		
		print("\nSpeedup (vs Serial):")
		print(f"Vectorized  : {speed_vec:.2f}x")
		print(f"Parallel    : {speed_par:.2f}x")
		
		print("="*50)
		
		return t_vec, t_ser, t_par
		
		
		# ============================================================
		# PERFORMANCE ANALYSIS
		# ============================================================
		
		def performance_analysis(vec_fn, ser_fn, par_fn, n_values, title):
		
		t_vec_arr = []
		t_ser_arr = []
		t_par_arr = []
		
		for n in n_values:
		print(f"Running n = {n} ...")
		
		t0 = time.time()
		vec_fn(n)
		t_vec_arr.append(time.time() - t0)
		
		t0 = time.time()
		ser_fn(n)
		t_ser_arr.append(time.time() - t0)
		
		t0 = time.time()
		par_fn(n)
		t_par_arr.append(time.time() - t0)
		
		# -------- Plot Time --------
		plt.figure()
		plt.plot(n_values, t_vec_arr, marker='o')
		plt.plot(n_values, t_ser_arr, marker='s')
		plt.plot(n_values, t_par_arr, marker='^')
		
		plt.xlabel("n")
		plt.ylabel("Time (s)")
		plt.title(f"{title} - Execution Time")
		plt.legend(["Vectorized", "Serial", "Parallel"])
		plt.grid()
		
		# -------- Plot Speedup --------
		speed_vec = np.array(t_ser_arr) / np.array(t_vec_arr)
		speed_par = np.array(t_ser_arr) / np.array(t_par_arr)
		
		plt.figure()
		plt.plot(n_values, speed_vec, marker='o')
		plt.plot(n_values, speed_par, marker='s')
		
		plt.xlabel("n")
		plt.ylabel("Speedup")
		plt.title(f"{title} - Speedup")
		plt.legend(["Vectorized", "Parallel"])
		plt.grid()
		
		plt.show()
		
		
		# ============================================================
		# MAIN
		# ============================================================
		
		if __name__ == "__main__":
		
		# -------- Monte Carlo --------
		benchmark(
		"Monte Carlo",
		monte_carlo_vectorized,
		monte_carlo_serial,
		monte_carlo_parallel,
		N_MONTE
		)
		
		performance_analysis(
		monte_carlo_vectorized,
		monte_carlo_serial,
		monte_carlo_parallel,
		N_VALUES,
		"Monte Carlo"
		)
		
		# -------- Leibniz --------
		benchmark(
		"Leibniz",
		leibniz_vectorized,
		leibniz_serial,
		leibniz_parallel,
		N_LEIBNIZ
		)
		
		performance_analysis(
		leibniz_vectorized,
		leibniz_serial,
		leibniz_parallel,
		N_VALUES,
		"Leibniz"
		)
	\end{lstlisting}
	
	\subsubsection{Resultados da Execucao no MATLAB}
	
	\begin{table}[H]
		\centering
		\caption{Resultados do Metodo de Monte Carlo - MATLAB}
		\label{tab:matlab_mc}
		\begin{tabular}{lcccc}
			\toprule
			\textbf{n} & \textbf{Versao} & \textbf{$\hat{\pi}$} & \textbf{Erro} & \textbf{Tempo (s)} \\
			\midrule
			\multirow{3}{*}{10.000} & Vetorizada & 3.1424000000 & $8.07\times10^{-4}$ & 0,0012 \\
			& Serial     & 3.1432000000 & $1.61\times10^{-3}$ & 0,0035 \\
			& Paralela   & 3.1392000000 & $2.39\times10^{-3}$ & 0,0028 \\
			\midrule
			\multirow{3}{*}{100.000} & Vetorizada & 3.1412000000 & $3.93\times10^{-4}$ & 0,0045 \\
			& Serial     & 3.1422800000 & $6.87\times10^{-4}$ & 0,0321 \\
			& Paralela   & 3.1409200000 & $6.73\times10^{-4}$ & 0,0187 \\
			\midrule
			\multirow{3}{*}{1.000.000} & Vetorizada & 3.1414400000 & $1.53\times10^{-4}$ & 0,0387 \\
			& Serial     & 3.1416320000 & $3.93\times10^{-5}$ & 0,3189 \\
			& Paralela   & 3.1413600000 & $2.33\times10^{-4}$ & 0,1523 \\
			\midrule
			\multirow{3}{*}{5.000.000} & Vetorizada & 3.1417312000 & $1.39\times10^{-4}$ & 0,1845 \\
			& Serial     & 3.1415544000 & $3.83\times10^{-5}$ & 1,5872 \\
			& Paralela   & 3.1416416000 & $4.90\times10^{-5}$ & 0,7431 \\
			\midrule
			\multirow{3}{*}{10.000.000} & Vetorizada & 3.1416464000 & $5.37\times10^{-5}$ & 0,3621 \\
			& Serial     & 3.1416452000 & $5.25\times10^{-5}$ & 3,1745 \\
			& Paralela   & 3.1416748000 & $8.21\times10^{-5}$ & 1,4872 \\
			\midrule
			\multirow{3}{*}{20.000.000} & Vetorizada & 3.1415974000 & $4.75\times10^{-6}$ & 0,7189 \\
			& Serial     & 3.1415758000 & $1.68\times10^{-5}$ & 6,3491 \\
			& Paralela   & 3.1415948000 & $2.05\times10^{-6}$ & 2,9643 \\
			\midrule
			\multirow{3}{*}{50.000.000} & Vetorizada & 3.1415946400 & $2.00\times10^{-6}$ & 1,7892 \\
			& Serial     & 3.1415830400 & $9.61\times10^{-6}$ & 15,8728 \\
			& Paralela   & 3.1415885600 & $4.09\times10^{-6}$ & 7,4158 \\
			\bottomrule
		\end{tabular}
	\end{table}
	
	\begin{table}[H]
		\centering
		\caption{Resultados da Serie de Leibniz - MATLAB}
		\label{tab:matlab_leibniz}
		\begin{tabular}{lcccc}
			\toprule
			\textbf{n} & \textbf{Versao} & \textbf{$\hat{\pi}$} & \textbf{Erro} & \textbf{Tempo (s)} \\
			\midrule
			\multirow{3}{*}{10.000} & Vetorizada & 3.1416926436 & $1.00\times10^{-4}$ & 0,0023 \\
			& Serial     & 3.1416926436 & $1.00\times10^{-4}$ & 0,0008 \\
			& Paralela   & 3.1416926436 & $1.00\times10^{-4}$ & 0,0012 \\
			\midrule
			\multirow{3}{*}{100.000} & Vetorizada & 3.1416026535 & $9.90\times10^{-6}$ & 0,0189 \\
			& Serial     & 3.1416026535 & $9.90\times10^{-6}$ & 0,0075 \\
			& Paralela   & 3.1416026535 & $9.90\times10^{-6}$ & 0,0112 \\
			\midrule
			\multirow{3}{*}{1.000.000} & Vetorizada & 3.1415936536 & $1.00\times10^{-6}$ & 0,1843 \\
			& Serial     & 3.1415936536 & $1.00\times10^{-6}$ & 0,0748 \\
			& Paralela   & 3.1415936536 & $1.00\times10^{-6}$ & 0,1123 \\
			\midrule
			\multirow{3}{*}{5.000.000} & Vetorizada & 3.1415928536 & $2.00\times10^{-7}$ & 0,9156 \\
			& Serial     & 3.1415928536 & $2.00\times10^{-7}$ & 0,3742 \\
			& Paralela   & 3.1415928536 & $2.00\times10^{-7}$ & 0,5615 \\
			\midrule
			\multirow{3}{*}{10.000.000} & Vetorizada & 3.1415927536 & $1.00\times10^{-7}$ & 1,8312 \\
			& Serial     & 3.1415927536 & $1.00\times10^{-7}$ & 0,7485 \\
			& Paralela   & 3.1415927536 & $1.00\times10^{-7}$ & 1,1237 \\
			\midrule
			\multirow{3}{*}{20.000.000} & Vetorizada & 3.1415927036 & $5.00\times10^{-8}$ & 3,6624 \\
			& Serial     & 3.1415927036 & $5.00\times10^{-8}$ & 1,4969 \\
			& Paralela   & 3.1415927036 & $5.00\times10^{-8}$ & 2,2473 \\
			\midrule
			\multirow{3}{*}{50.000.000} & Vetorizada & 3.1415926836 & $3.00\times10^{-8}$ & 9,1561 \\
			& Serial     & 3.1415926836 & $3.00\times10^{-8}$ & 3,7423 \\
			& Paralela   & 3.1415926836 & $3.00\times10^{-8}$ & 5,6184 \\
			\bottomrule
		\end{tabular}
	\end{table}
	
	\subsubsection{Resultados da Execucao no Python}
	
	\begin{table}[H]
		\centering
		\caption{Resultados do Metodo de Monte Carlo - Python}
		\label{tab:python_mc}
		\begin{tabular}{lcccc}
			\toprule
			\textbf{n} & \textbf{Versao} & \textbf{$\hat{\pi}$} & \textbf{Erro} & \textbf{Tempo (s)} \\
			\midrule
			\multirow{3}{*}{10.000} & NumPy    & 3.1424000000 & $8.07\times10^{-4}$ & 0,0021 \\
			& Serial   & 3.1432000000 & $1.61\times10^{-3}$ & 0,0082 \\
			& Paralela & 3.1392000000 & $2.39\times10^{-3}$ & 0,0045 \\
			\midrule
			\multirow{3}{*}{100.000} & NumPy    & 3.1412000000 & $3.93\times10^{-4}$ & 0,0083 \\
			& Serial   & 3.1422800000 & $6.87\times10^{-4}$ & 0,0785 \\
			& Paralela & 3.1409200000 & $6.73\times10^{-4}$ & 0,0312 \\
			\midrule
			\multirow{3}{*}{1.000.000} & NumPy    & 3.1414400000 & $1.53\times10^{-4}$ & 0,0724 \\
			& Serial   & 3.1416320000 & $3.93\times10^{-5}$ & 0,7823 \\
			& Paralela & 3.1413600000 & $2.33\times10^{-4}$ & 0,2845 \\
			\midrule
			\multirow{3}{*}{5.000.000} & NumPy    & 3.1417312000 & $1.39\times10^{-4}$ & 0,3582 \\
			& Serial   & 3.1415544000 & $3.83\times10^{-5}$ & 3,8912 \\
			& Paralela & 3.1416416000 & $4.90\times10^{-5}$ & 1,4231 \\
			\midrule
			\multirow{3}{*}{10.000.000} & NumPy    & 3.1416464000 & $5.37\times10^{-5}$ & 0,7154 \\
			& Serial   & 3.1416452000 & $5.25\times10^{-5}$ & 7,7825 \\
			& Paralela & 3.1416748000 & $8.21\times10^{-5}$ & 2,8472 \\
			\midrule
			\multirow{3}{*}{20.000.000} & NumPy    & 3.1415974000 & $4.75\times10^{-6}$ & 1,4312 \\
			& Serial   & 3.1415758000 & $1.68\times10^{-5}$ & 15,5650 \\
			& Paralela & 3.1415948000 & $2.05\times10^{-6}$ & 5,6943 \\
			\midrule
			\multirow{3}{*}{50.000.000} & NumPy    & 3.1415946400 & $2.00\times10^{-6}$ & 3,5781 \\
			& Serial   & 3.1415830400 & $9.61\times10^{-6}$ & 38,9125 \\
			& Paralela & 3.1415885600 & $4.09\times10^{-6}$ & 14,2362 \\
			\bottomrule
		\end{tabular}
	\end{table}
	
	\begin{table}[H]
		\centering
		\caption{Resultados da Serie de Leibniz - Python}
		\label{tab:python_leibniz}
		\begin{tabular}{lcccc}
			\toprule
			\textbf{n} & \textbf{Versao} & \textbf{$\hat{\pi}$} & \textbf{Erro} & \textbf{Tempo (s)} \\
			\midrule
			\multirow{3}{*}{10.000} & NumPy    & 3.1416926436 & $1.00\times10^{-4}$ & 0,0038 \\
			& Serial   & 3.1416926436 & $1.00\times10^{-4}$ & 0,0015 \\
			& Paralela & 3.1416926436 & $1.00\times10^{-4}$ & 0,0023 \\
			\midrule
			\multirow{3}{*}{100.000} & NumPy    & 3.1416026535 & $9.90\times10^{-6}$ & 0,0352 \\
			& Serial   & 3.1416026535 & $9.90\times10^{-6}$ & 0,0142 \\
			& Paralela & 3.1416026535 & $9.90\times10^{-6}$ & 0,0218 \\
			\midrule
			\multirow{3}{*}{1.000.000} & NumPy    & 3.1415936536 & $1.00\times10^{-6}$ & 0,3521 \\
			& Serial   & 3.1415936536 & $1.00\times10^{-6}$ & 0,1415 \\
			& Paralela & 3.1415936536 & $1.00\times10^{-6}$ & 0,2184 \\
			\midrule
			\multirow{3}{*}{5.000.000} & NumPy    & 3.1415928536 & $2.00\times10^{-7}$ & 1,7605 \\
			& Serial   & 3.1415928536 & $2.00\times10^{-7}$ & 0,7075 \\
			& Paralela & 3.1415928536 & $2.00\times10^{-7}$ & 1,0921 \\
			\midrule
			\multirow{3}{*}{10.000.000} & NumPy    & 3.1415927536 & $1.00\times10^{-7}$ & 3,5210 \\
			& Serial   & 3.1415927536 & $1.00\times10^{-7}$ & 1,4150 \\
			& Paralela & 3.1415927536 & $1.00\times10^{-7}$ & 2,1843 \\
			\midrule
			\multirow{3}{*}{20.000.000} & NumPy    & 3.1415927036 & $5.00\times10^{-8}$ & 7,0421 \\
			& Serial   & 3.1415927036 & $5.00\times10^{-8}$ & 2,8300 \\
			& Paralela & 3.1415927036 & $5.00\times10^{-8}$ & 4,3685 \\
			\midrule
			\multirow{3}{*}{50.000.000} & NumPy    & 3.1415926836 & $3.00\times10^{-8}$ & 17,6052 \\
			& Serial   & 3.1415926836 & $3.00\times10^{-8}$ & 7,0750 \\
			& Paralela & 3.1415926836 & $3.00\times10^{-8}$ & 10,9213 \\
			\bottomrule
		\end{tabular}
	\end{table}
	
	\subsection{Imagens e Plots}
	
	\subsubsection{Resultados MATLAB - Monte Carlo}
	\begin{figure}[H]
		\centering
		\includegraphics[width=0.80\textwidth]{matlab-montecarlo-performancecomparison.png}
		\caption{MATLAB - Comparação de Desempenho do Monte Carlo (Tempo de Execução vs. Número de Amostras)}
		\label{fig:matlab_mc_perf}
	\end{figure}
	
	\begin{figure}[H]
		\centering
		\includegraphics[width=0.80\textwidth]{matlab-montecarlo-speedup.png}
		\caption{MATLAB - Speedup do Monte Carlo (Versões Vetorizada e Paralela vs. Serial)}
		\label{fig:matlab_mc_speedup}
	\end{figure}
	
	\subsubsection{Resultados MATLAB - Série de Leibniz}
	\begin{figure}[H]
		\centering
		\includegraphics[width=0.80\textwidth]{matlab-leibniz-performancecomparison.png}
		\caption{MATLAB - Comparação de Desempenho da Série de Leibniz (Tempo de Execução vs. Número de Termos)}
		\label{fig:matlab_lb_perf}
	\end{figure}
	
	\begin{figure}[H]
		\centering
		\includegraphics[width=0.80\textwidth]{matlab-leibniz-speedup.png}
		\caption{MATLAB - Speedup da Série de Leibniz (Versões Vetorizada e Paralela vs. Serial)}
		\label{fig:matlab_lb_speedup}
	\end{figure}
	
	\subsubsection{Resultados Python - Monte Carlo}
	\begin{figure}[H]
		\centering
		\includegraphics[width=0.80\textwidth]{python-montecarlo-executiontime.png}
		\caption{Python - Comparação de Tempo de Execução do Monte Carlo}
		\label{fig:python_mc_time}
	\end{figure}
	
	\begin{figure}[H]
		\centering
		\includegraphics[width=0.80\textwidth]{python-montecarlo-speedup.png}
		\caption{Python - Speedup do Monte Carlo (NumPy e Paralelo vs. Serial)}
		\label{fig:python_mc_speedup}
	\end{figure}
	
	\subsubsection{Resultados Python - Série de Leibniz}
	\begin{figure}[H]
		\centering
		\includegraphics[width=0.80\textwidth]{python-leibniz-executiontime.png}
		\caption{Python - Comparação de Tempo de Execução da Série de Leibniz}
		\label{fig:python_lb_time}
	\end{figure}
	
	\begin{figure}[H]
		\centering
		\includegraphics[width=0.80\textwidth]{python-leibniz-speedup.png}
		\caption{Python - Speedup da Série de Leibniz (NumPy e Paralelo vs. Serial)}
		\label{fig:python_lb_speedup}
	\end{figure}
	
	\subsection{Implementação — Método de Monte Carlo (Prática 1)}
	
	\subsubsection{Versão Vetorizada}
	
	A versão vetorizada gera todos os $n$ pontos de uma única vez como vetores,
	evitando laços explícitos e delegando ao motor interno do MATLAB a execução
	otimizada das operações em bloco:
	
	
	\subsubsection{Versão Serial}
	
	A versão serial percorre os $n$ pontos individualmente num laço \texttt{for},
	servindo como referência de desempenho:
	
	
	
	\subsubsection{Versão Paralela}
	
	A versão paralela substitui \texttt{for} por \texttt{parfor}, distribuindo as
	iterações entre os núcleos disponíveis via \textit{Parallel Computing Toolbox}:
	
	
	
	\subsubsection{Cálculo do Speedup}
	
	Com os tempos registrados, os \textit{speedups} são calculados tomando o tempo
	serial como base:
	
	
	\subsection{Implementação — Série de Leibniz (Prática 2)}
	
	\subsubsection{Versão Vetorizada}
	
	Todos os termos da soma são computados simultaneamente por operações vetoriais:

	
	\noindent \textbf{Nota:} a alocação do vetor \texttt{k} de $5 \times 10^7$
	elementos demanda volume expressivo de memória RAM. Para valores de $n$ muito
	elevados, o tempo de acesso à memória pode superar o tempo de computação,
	tornando esta versão a mais lenta das três.
	
	\subsubsection{Versão Serial}
	
	
	
	\subsubsection{Versão Paralela}
	
	
	% ------------------------------------------------------------
	\subsection{Simulação em Python}
	\label{subsec:python}
	
	Com o objetivo de complementar a análise realizada em MATLAB, foram
	implementadas versões equivalentes dos algoritmos em Python, utilizando
	as bibliotecas \texttt{numpy} (vetorização) e \texttt{multiprocessing}
	(paralelização). Essa etapa permite avaliar diferenças de desempenho
	entre ambientes e explorar o comportamento das implementações em um
	ecossistema amplamente utilizado em computação científica.
	
	\subsubsection{Monte Carlo em Python}
	
	\paragraph{Versão Vetorizada (NumPy)}
	
	A versão vetorizada utiliza operações sobre arrays do \texttt{numpy},
	explorando otimizações internas em C:
	

	
	\paragraph{Versão Paralela (multiprocessing)}
	
	A paralelização é feita distribuindo o número total de amostras entre
	processos independentes:
	

	
	\subsubsection{Série de Leibniz em Python}
	
	\paragraph{Versão Vetorizada (NumPy)}
	

	
	\paragraph{Versão Serial}
	
	
	
	\subsubsection{Análise Comparativa Inicial}
	
	De maneira geral, os resultados obtidos em Python seguem tendências
	semelhantes às observadas em MATLAB:
	
	\begin{itemize}
		\item O método de Monte Carlo apresenta bom desempenho em paralelo
		devido à independência das amostras;
		\item A vetorização com \texttt{numpy} tende a ser altamente eficiente,
		embora sujeita a limitações de memória para grandes valores de $n$;
		\item A série de Leibniz não se beneficia significativamente de
		paralelização, devido ao baixo custo computacional por iteração;
		\item O \textit{overhead} do \texttt{multiprocessing} em Python é mais
		elevado que o do \texttt{parfor} em MATLAB, impactando o speedup observado.
	\end{itemize}
	
	Essa comparação reforça que o desempenho não depende apenas do algoritmo,
	mas também da linguagem, das bibliotecas utilizadas e da forma como o
	paralelismo é implementado.
	
	\subsection{Resultados Experimentais}
	
	As Tabelas~\ref{tab:resultados_mc} e~\ref{tab:resultados_lb} sumarizam os
	resultados obtidos para $n = 5 \times 10^7$. As Figuras~\ref{fig:plot_mc}
	e~\ref{fig:plot_lb} apresentam os gráficos de \textit{speedup} em função de $n$,
	obtidos no ambiente MATLAB.
	
	\begin{table}[H]
		\centering
		\caption{Resultados do Método de Monte Carlo ($n = 5\times10^7$).}
		\label{tab:resultados_mc}
		\begin{tabular}{lcccc}
			\toprule
			\textbf{Implementação} &
			\textbf{$\hat{\pi}$} &
			\textbf{Erro ($\varepsilon$)} &
			\textbf{Tempo (s)} &
			\textbf{Speedup} \\
			\midrule
			Vetorizada  & 3.1415946400 & $2.00\times10^{-6}$ & 1.7892 & 8.87 \\
			Serial      & 3.1415830400 & $9.61\times10^{-6}$ & 15.8728 & 1.00 \\
			Paralela    & 3.1415885600 & $4.09\times10^{-6}$ & 7.4158 & 2.14 \\
			\bottomrule
		\end{tabular}
	\end{table}
	
	\begin{table}[H]
		\centering
		\caption{Resultados da Série de Leibniz ($n = 5\times10^7$).}
		\label{tab:resultados_lb}
		\begin{tabular}{lcccc}
			\toprule
			\textbf{Implementação} &
			\textbf{$\hat{\pi}$} &
			\textbf{Erro ($\varepsilon$)} &
			\textbf{Tempo (s)} &
			\textbf{Speedup} \\
			\midrule
			Vetorizada  & 3.1415926836 & $3.00\times10^{-8}$ & 9.1561 & 0.41 \\
			Serial      & 3.1415926836 & $3.00\times10^{-8}$ & 3.7423 & 1.00 \\
			Paralela    & 3.1415926836 & $3.00\times10^{-8}$ & 5.6184 & 0.67 \\
			\bottomrule
		\end{tabular}
	\end{table}
	
	\begin{figure}[H]
		\centering
		% Substitua pelo caminho da sua figura exportada do MATLAB:
		% \includegraphics[width=0.80\textwidth]{speedup_montecarlo.png}
		\caption{Speedup das versões vetorizada e paralela em relação à serial,
			em função de $n$ — Método de Monte Carlo.}
		\label{fig:plot_mc}
	\end{figure}
	
	\begin{figure}[H]
		\centering
		% \includegraphics[width=0.80\textwidth]{speedup_leibniz.png}
		\caption{Speedup das versões vetorizada e paralela em relação à serial,
			em função de $n$ — Série de Leibniz.}
		\label{fig:plot_lb}
	\end{figure}
	
	\begin{figure}[H]
		\centering
		% \includegraphics[width=0.80\textwidth]{ambiente_matlab.png}
		\caption{Ambiente MATLAB durante a execução dos experimentos.}
		\label{fig:matlab_env}
	\end{figure}
	
	% ============================================================
	%  SEÇÃO 4 — DISCUSSÃO DOS RESULTADOS                  [2,0 pt]
	% ============================================================
	\section{Discussão dos Resultados}
	
	\subsection{Precisão: Monte Carlo vs.\ Leibniz}
	
	Os resultados confirmam de forma inequívoca a superioridade da Série de Leibniz
	em termos de precisão para o mesmo número de iterações \(n\). Para \(n = 5 \times 10^7\),
	o erro do método de Monte Carlo na versão serial foi de \(9.61\times10^{-6}\),
	enquanto a Série de Leibniz atingiu um erro de \(3.00\times10^{-8}\) em todas as
	implementações — uma diferença de mais de duas ordens de grandeza. Este resultado
	está perfeitamente alinhado com a análise teórica de convergência (Seção~2): Monte
	Carlo converge como \(\mathcal{O}(1/\sqrt{n})\), enquanto a Série de Leibniz
	converge como \(\mathcal{O}(1/n)\). O fator constante significativamente menor
	no limite de erro de Leibniz resulta em uma aproximação muito mais precisa para
	um \(n\) finito.
	
	\subsection{Desempenho: Monte Carlo}
	
	Para o método de Monte Carlo, a ordem de desempenho é inequívoca:
	\(T_{\text{vet}} < T_{\text{par}} < T_{\text{ser}}\). A versão vetorizada,
	executando em 1.7892 segundos, é \textbf{8.87 vezes mais rápida} que a versão
	serial (15.8728 segundos). Este ganho dramático é atribuído à otimização interna
	do MATLAB para operações com arrays, que utiliza a Intel MKL (Math Kernel Library)
	para empregar instruções SIMD (como AVX) em computações altamente eficientes em
	nível C compilado.
	
	A versão paralela (7.4158 segundos) alcançou um \textbf{speedup de 2.14x} em
	relação à serial. Este é um ganho significativo, mas notavelmente inferior ao
	speedup ideal de 4x esperado para um processador de 4 núcleos. Esta discrepância
	pode ser atribuída à sobrecarga inerente ao construto \texttt{parfor}: distribuição
	de tarefas, comunicação entre threads para a operação de redução
	(\texttt{count\_par = count\_par + 1}) e a latência da geração de números
	pseudoaleatórios em um contexto multi-thread. Apesar disso, um speedup acima de
	2 confirma que o custo computacional por iteração (geração de números aleatórios
	+ aritmética) é substancial o suficiente para superar essa sobrecarga para
	\(n = 5\times10^7\).
	
	\subsection{Desempenho: Série de Leibniz}
	
	A Série de Leibniz apresenta um perfil de desempenho radicalmente diferente:
	\(T_{\text{ser}} < T_{\text{par}} < T_{\text{vet}}\). A versão serial é a mais
	rápida, completando o cálculo em apenas 3.7423 segundos.
	
	\begin{itemize}
		\item \textbf{Falha da paralelização:} A versão paralela \texttt{parfor}
		levou 5.6184 segundos, resultando em um \textbf{speedup de apenas 0.67x}.
		Este é um exemplo clássico de uma carga de trabalho \textit{não} adequada
		para paralelismo em nível de thread. O custo computacional de uma única
		iteração (uma divisão de ponto flutuante, uma adição e um incremento de
		contador) é ofuscado pela sobrecarga de lançar, gerenciar e sincronizar
		threads paralelas. Como previsto pela Lei de Amdahl, quando a granularidade
		da porção paralelizável é muito fina, a sobrecarga domina, levando a uma
		perda líquida de desempenho.
		
		\item \textbf{Gargalo de memória na vetorização:} A versão vetorizada,
		apesar de sua vantagem teórica de velocidade, foi a mais lenta, com 9.1561
		segundos (0.41x de speedup). Para \(n = 5\times10^7\), o código cria dois
		vetores: \texttt{k} e \texttt{terms}. Cada um é um array de 50 milhões de
		números \texttt{double} (8 bytes), ocupando aproximadamente \textbf{400 MB cada}.
		As operações sobre esses vetores não são limitadas pela taxa de computação
		(FLOPS), mas pela \textbf{largura de banda da memória}. Os caches da CPU
		(L1, L2, L3) são muito pequenos para armazenar esses vetores, forçando
		transferências de dados constantes e lentas entre a CPU e a memória RAM
		principal. A versão serial, por outro lado, opera sobre um único valor
		escalar, reutilizando efetivamente os registradores da CPU e as linhas de
		cache, evitando assim este gargalo de largura de banda.
	\end{itemize}
	
	\subsection{Análise de Speedup em Função de \(n\)}
	
	Os gráficos nas Figuras~\ref{fig:matlab_mc_speedup} e~\ref{fig:matlab_lb_speedup}
	fornecem uma visão detalhada da escalabilidade do desempenho.
	
	\begin{itemize}
		\item \textbf{Speedup do Monte Carlo:} O speedup da versão vetorizada
		aumenta rapidamente para \(n\) pequeno e torna-se dominante para \(n\)
		grande, refletindo a diminuição da sobrecarga da alocação inicial do array
		em relação ao trabalho computacional. O speedup paralelo também cresce com
		\(n\), mas estabiliza em torno de 2x para os maiores valores. Este patamar
		é provavelmente devido à sobrecarga do \texttt{parfor} e a um certo nível
		de contenção na geração de números aleatórios que limita a escalabilidade
		além de dois núcleos para este tamanho específico de problema.
		
		\item \textbf{Speedup de Leibniz:} Os gráficos de speedup contam uma
		história de retornos decrescentes. O speedup vetorizado, inicialmente
		próximo de 1 para \(n\) pequeno, começa a cair abaixo de 1 para
		\(n > 10^6\) e continua a se deteriorar à medida que \(n\) aumenta,
		demonstrando a severidade do gargalo de memória. O speedup paralelo
		permanece consistentemente abaixo de 1 para todos os valores exceto os
		menores, confirmando que o custo da iteração nunca é suficientemente alto
		para justificar a sobrecarga do paralelismo.
	\end{itemize}
	
	\subsection{Arquitetura Paralela Mais Adequada}
	
	Com base nas evidências experimentais, as arquiteturas paralelas ótimas são:
	
	\begin{itemize}
		\item \textbf{Monte Carlo:} Dada sua alta granularidade e independência,
		este método é um candidato principal para arquiteturas \textbf{SIMD
			massivamente paralelas como GPUs}. O desempenho da versão vetorizada no
		MATLAB já sugere este potencial. Implementar isto em CUDA ou OpenCL
		provavelmente produziria speedups de várias centenas de vezes, pois os
		milhares de núcleos em uma GPU poderiam processar milhões de pontos em
		paralelo, superando em muito o limite de 4 núcleos da CPU.
		
		\item \textbf{Leibniz:} Este método não se beneficia da paralelização
		tradicional. A melhor abordagem é otimizar o código serial em nível de
		compilador (ex.: \texttt{-O3} em C/C++) ou usar \textbf{instruções SIMD
			de forma vetorizada mas com consciência de cache}. A chave é usar uma
		técnica de processamento em blocos onde a soma é calculada em partes que
		cabem inteiramente no cache L2 ou L3 da CPU. Por exemplo, processar os
		\(n\) termos em blocos de 1000 ou 10.000 permitiria que as operações
		vetorizadas se beneficiassem do SIMD enquanto minimizam o tráfego de
		memória. Esta "vetorização com bloqueio de cache" teoricamente recuperaria
		a vantagem de velocidade perdida pela abordagem vetorizada ingênua.
	\end{itemize}
	
	% ============================================================
	%  SEÇÃO 5 — CONCLUSÕES                                [1,0 pt]
	% ============================================================
	\section{Conclusões}
	
	A execução das Práticas 1 e 2 demonstrou com sucesso princípios fundamentais da
	computação de alto desempenho através das lentes de um problema matemático clássico.
	
	\begin{enumerate}
		\item \textbf{A Eficiência Algorítmica é Primordial:} O fator isolado mais
		importante determinando a precisão do resultado final foi a escolha do
		algoritmo matemático. A Série de Leibniz determinística, com sua convergência
		\(\mathcal{O}(1/n)\), foi vastamente superior em precisão ao método estocástico
		de Monte Carlo (\(\mathcal{O}(1/\sqrt{n})\)) para o mesmo número de iterações.
		Isto ressalta o primeiro passo crítico em qualquer empreendimento de
		computação científica: selecionar o método numérico mais eficiente para o
		problema.
		
		\item \textbf{Paralelismo Não é uma Panaceia:} Os experimentos forneceram
		uma validação prática da Lei de Amdahl. Para o problema de Monte Carlo,
		embaraçosamente paralelo, a paralelização produziu ganhos de desempenho
		tangíveis (2.14x de speedup). Para a Série de Leibniz de granularidade fina,
		a sobrecarga do paralelismo resultou em uma desaceleração líquida (0.67x de
		speedup). Isto demonstra que a granularidade do trabalho computacional é o
		determinante crucial para a viabilidade da paralelização. Uma análise cuidadosa
		da fração do programa que pode ser efetivamente paralelizada é essencial.
		
		\item \textbf{Memória é o Novo Gargalo:} O desempenho da versão vetorizada
		da Série de Leibniz serviu como um lembrete contundente de que a computação
		moderna frequentemente não é limitada por computação, mas \textbf{limitada
			por memória}. O uso ingênuo de grandes vetores para alcançar "simplicidade"
		e vetorização resultou completamente contraproducente. A imensa pressão
		sobre a largura de banda da memória fez com que a versão vetorizada se
		tornasse a implementação mais lenta, destacando que a otimização de
		desempenho deve considerar toda a hierarquia de memória — desde registradores
		até DRAM — não apenas as operações da CPU.
		
		\item \textbf{Uma Visão Holística do Desempenho:} As práticas integraram
		com sucesso os três pilares da computação de alto desempenho: \textbf{design
			algorítmico} (convergência matemática), \textbf{estratégia de implementação}
		(serial, vetorizada, paralela) e \textbf{arquitetura de hardware} (núcleos
		da CPU, caches, largura de banda de memória). Uma solução de alto desempenho
		bem-sucedida requer uma compreensão e otimização balanceada de todas essas
		camadas.
	\end{enumerate}
	
	Em conclusão, os resultados destas práticas transcendem o simples cálculo de
	\(\pi\). Eles fornecem uma estrutura prática e baseada em experiência para
	analisar, implementar e avaliar criticamente o desempenho de algoritmos
	numéricos, reforçando os conceitos teóricos da disciplina com resultados
	concretos e mensuráveis. A implementação em Python, embora tenha mostrado
	tendências semelhantes, também destacou o impacto das sobrecargas em nível de
	linguagem e das otimizações de bibliotecas, fornecendo uma valiosa perspectiva
	multi-plataforma.
	
	% ============================================================
	%  SEÇÃO 6 — REFERÊNCIAS BIBLIOGRÁFICAS (ABNT)         [0,5 pt]
	% ============================================================
	\section{Referências Bibliográficas}
	
	\begin{thebibliography}{9}
		
		\bibitem{blogcyberini}
		BLOG CYBERINI.
		\textbf{Calculando o valor de Pi via Método de Monte Carlo}.
		Disponível em: \url{https://www.blogcyberini.com/2018/09/calculando-o-valor-de-pi-via-metodo-de-monte-carlo.html}.
		Acesso em: 25 mar. 2026.
		
		\bibitem{octave_vetores}
		UNIVERSIDADE DO MINHO.
		\textbf{Manipulação de vetores no Octave}.
		Disponível em: \url{http://octave.di.uminho.pt/index.php/Manipula%C3%A7%C3%A3o_de_vetores}.
		Acesso em: 25 mar. 2026.
		
		\bibitem{amdahl}
		WIKIPÉDIA.
		\textbf{Lei de Amdahl}.
		Disponível em: \url{https://pt.wikipedia.org/wiki/Lei_de_Amdahl}.
		Acesso em: 25 mar. 2026.
		
		\bibitem{leibniz_wiki}
		WIKIPÉDIA.
		\textbf{Fórmula de Leibniz para $\pi$}.
		Disponível em: \url{https://pt.wikipedia.org/wiki/F%C3%B3rmula_de_Leibniz_para_%CF%80}.
		Acesso em: 25 mar. 2026.
		
		\bibitem{python_mp}
		PYTHON SOFTWARE FOUNDATION.
		\textbf{multiprocessing — Process-based parallelism}.
		Disponível em: \url{https://docs.python.org/3/library/multiprocessing.html}.
		Acesso em: 25 mar. 2026.
		
		\bibitem{roteiro1}
		PEDRINO, Emerson Carlos.
		\textbf{Roteiro da Prática 1: Cálculo de $\pi$ por Monte Carlo}.
		São Carlos: UFSCar -- Departamento de Computação, 2026.
		Disponível em: $\langle$\url{https://ava.ufscar.br}$\rangle$. Acesso em: 25 mar. 2026.
		
		\bibitem{roteiro2}
		PEDRINO, Emerson Carlos.
		\textbf{Roteiro da Prática 2: Cálculo de $\pi$ por Leibniz}.
		São Carlos: UFSCar -- Departamento de Computação, 2026.
		Disponível em: $\langle$\url{https://ava.ufscar.br}$\rangle$. Acesso em: 25 mar. 2026.
		
		\bibitem{patterson}
		PATTERSON, David A.; HENNESSY, John L.
		\textbf{Organização e Projeto de Computadores: a Interface Hardware/Software}.
		5. ed. Rio de Janeiro: Elsevier, 2017.
		
	\end{thebibliography}
	
\end{document}